\setcounter{section}{0}

\Needspace{5\baselineskip}
\hypertarget{overview-of-world-models}{%
\section{Overview of World Models}\label{overview-of-world-models}}

\Needspace{5\baselineskip}
\hypertarget{i.-definition-of-a-world-model}{%
\subsection{I. Definition of a World Model}\label{i.-definition-of-a-world-model}}

LLMs already possess strong capabilities in the digital world, but their understanding of the physical world remains insufficient. World models offer a possible solution. If the core of an LLM is ``next-token prediction,'' the core of a world model is ``next-state prediction.'' A world model is defined as a system that predicts the next state given the current state and action (in practice, this often means several future states; the definition is simplified here):

P(st+1\textbar st,at)

As a model of the physical world, a world model should handle multimodal inputs (not limited to vision, but including sensor data, spectral data, and other modalities), generate explicit or implicit representations of world states, and effectively predict the next state under a given action.

\Needspace{5\baselineskip}
\hypertarget{ii.-the-origins-of-world-models-world-models-in-rl}{%
\subsection{II. The Origins of World Models: World Models in RL}\label{ii.-the-origins-of-world-models-world-models-in-rl}}

\Needspace{5\baselineskip}
\hypertarget{core-idea-5}{%
\subsubsection{(1) Core Idea}\label{core-idea-5}}

When humans perceive the world, we construct a mental model from our limited sensory information. Our decisions and actions are based on this internal model. Inspired by this idea from cognitive science, the paper explores building generative neural-network models for reinforcement-learning (RL) environments.

The researchers divide the agent into a world model and a control model. The world model learns a compressed representation of the environment in space and time; the control model uses features extracted from the world model as inputs to train a task-solving policy.

The advantage of this approach is that it not only handles credit assignment efficiently, but also allows the agent to train itself in virtual ``dreams'' generated by its world model before transferring the policy to the real environment. This is model-based reinforcement learning.

\Needspace{5\baselineskip}
\hypertarget{world-model-architecture-and-training}{%
\subsubsection{(2) World-Model Architecture and Training}\label{world-model-architecture-and-training}}

\begin{enumerate}
\def\labelenumi{\arabic{enumi}.}
\item
  Perception model: Given multimodal information such as visual pixels at time t, use a variational autoencoder (VAE) or a similar model to compress it into a small vector, the latent z\_t.
\item
  Prediction model: Maintain hidden state h\_t and record historical memory with an RNN-like structure to ensure consistency with physical laws over long periods. Given hidden state h\_t, current-state information z\_t, and action a\_t, output the future-state prediction z*\_t+1. The model can directly output a sample or a probability-density function over future states. Learning uses a loss based on the discrepancy between predicted z*\_t+1 and the latent z\_t+1 corresponding to the true next state.
\end{enumerate}

\Needspace{5\baselineskip}
\hypertarget{collaboration-between-the-control-model-and-world-model}{%
\subsubsection{(3) Collaboration between the Control Model and World Model}\label{collaboration-between-the-control-model-and-world-model}}

\Needspace{5\baselineskip}
The control model consists of policy networks, value networks, and similar components trained through reinforcement learning. There are generally three ways for it to collaborate with the world model:

\begin{enumerate}
\def\labelenumi{\arabic{enumi}.}
\item
  The control model's inputs do not include world-model outputs, but those outputs are incorporated when performing operations such as group sampling of the control model's outputs.
\item
  Feed world-model outputs into the control model as well.
\item
  Merge the two directly, retaining multiple output heads.
\end{enumerate}

\Needspace{5\baselineskip}
\hypertarget{training-workflow}{%
\subsubsection{(4) Training Workflow}\label{training-workflow}}

\Needspace{5\baselineskip}
Training repeats the following cycle:

\begin{enumerate}
\def\labelenumi{\arabic{enumi}.}
\item
  Freeze model weights while the agent interacts with the real environment to collect state--action pairs.
\item
  Train the perception model on these data.
\item
  Freeze the perception model and train the prediction model.
\item
  Finally, freeze the perception and prediction models and train the control model.
\end{enumerate}

\Needspace{5\baselineskip}
\hypertarget{adversarial-policies-and-their-prevention}{%
\subsubsection{(5) Adversarial Policies and Their Prevention}\label{adversarial-policies-and-their-prevention}}

Because the world model is only an approximate probabilistic model of the real environment, the controller may discover and exploit its imperfections during optimization, producing ``adversarial policies'' (for example, moving in a particular way that causes generated fireballs to disappear). Appropriately increasing the prediction model's temperature raises uncertainty and stochasticity in the virtual environment, effectively preventing exploitation of these model defects. This helps the agent learn more robust policies and perform better in the real environment.

\Needspace{5\baselineskip}
\hypertarget{iii.-the-relationship-between-world-models-and-embodied-ai}{%
\subsection{III. The Relationship between World Models and Embodied AI}\label{iii.-the-relationship-between-world-models-and-embodied-ai}}

\Needspace{5\baselineskip}
\hypertarget{world-models-require-embodied-interaction}{%
\subsubsection{(1) World Models Require Embodied Interaction}\label{world-models-require-embodied-interaction}}

Watching videos alone is insufficient for a model to truly understand physical-world laws. Video data has several significant problems.

\begin{enumerate}
\def\labelenumi{\arabic{enumi}.}
\item
  Lack of action conditioning: Videos contain state changes but no action conditions. Genuine physical understanding often comes from interventions: ``If I push it, will it fall?'' ``If I heat it, will it change color?'' These often depend on action-conditioned data. Without action conditioning, one is simply training a video-generation model.
\item
  Lack of tactile and mechanical data: Current world-model training remains primarily visual. However, to genuinely understand physical laws, a world model must not only ``see and understand'' the world but also understand physical properties such as friction and hardness. This requires tactile and mechanical data to avoid ``visually plausible but physically incorrect'' predictions, such as objects without supporting forces or collisions without momentum transfer.
\end{enumerate}

\Needspace{5\baselineskip}
\hypertarget{world-models-are-crucial-to-the-development-of-embodied-ai}{%
\subsubsection{(2) World Models Are Crucial to the Development of Embodied AI}\label{world-models-are-crucial-to-the-development-of-embodied-ai}}

\begin{enumerate}
\def\labelenumi{\arabic{enumi}.}
\item
  A world model is an important component of a world--action model, combining with a policy model to form an embodied foundation model.
\item
  World models can serve as training and evaluation environments. Real-world trial and error is highly inefficient and can damage expensive hardware. Embodied agents with world models can operate in internally generated spaces, using those models as training and evaluation environments. Several points require attention in practice:
\end{enumerate}

\begin{enumerate}
\def\labelenumi{(\arabic{enumi})}
\item
  Because predicting world states (video generation) consumes substantially more compute than predicting actions, we need not predict frame by frame and can instead use chunk-level prediction.
\item
  Robots often have multiple cameras, whose images must be consistent in three dimensions, jointly forming a coherent 3D world. One approach first generates one view, places that image in the context, and uses it to condition generation of the next view, continuing in this way.
\end{enumerate}

\begin{enumerate}
\def\labelenumi{\arabic{enumi}.}
\setcounter{enumi}{2}
\tightlist
\item
  World models can serve as synthetic-data engines. They can generate large numbers of synthetic trajectories with interaction logic involving physical feedback, supporting action-policy training at very low marginal cost.
\end{enumerate}

\FloatBarrier
\Needspace{5\baselineskip}
\hypertarget{references-154}{%
\subsection{References}\label{references-154}}

\vspace{0.25em}
\begin{itemize}
\tightlist
\item
  Ha, D., \& Schmidhuber, J. (2018). \href{https://arxiv.org/abs/1803.10122}{World Models}. arXiv:1803.10122.
\item
  Sutton, R. S. (1991). \href{https://dl.acm.org/doi/10.1145/122344.122377}{Dyna, an Integrated Architecture for Learning, Planning, and Reacting}. ACM SIGART Bulletin.
\end{itemize}

\clearpage
